\documentclass{article}
\usepackage{amsmath}
\usepackage{amssymb}

\title{核心数学模型}
\date{}

\usepackage[UTF8]{ctex}
\begin{document}

\maketitle

\section*{受燃料限制的单位承诺}

为了使该模型能够直接在 Gurobi 中实施，我们使用了明确的初始条件$(y_0,x_0)$，以及实例数据中提供的、与特定时期相对应的分段线性生产成本曲线。所有数字都保持不变。 具体而言，对每个时段 $t$，求解器从 \texttt{f\_t.breakpoints\_MW} 和 \texttt{f\_t.values} 读取按顺序排列的断点和值的配对 $(\chi_{tm},\phi_{tm})_{m\in\mathcal{M}_t}$，然后通过 Gurobi 的 \texttt{addGenConstrPWL} 施加相应的分段线性约束。 当这些数据定义了一个凸曲线时，这就相当于一条支撑线上的曲线；否则，所实现的模型就会遵循非凸多段线本身的特征。 存储库的实现还提供了一个表示周期长度的参数$\Delta t = \texttt{period\_length}$；在当前的求解器中，$\Delta t$仅用于缩放生产成本项，而收入项$\xi_t x_t$和燃料预算$Q$则保持未缩放状态。

\begin{equation}
\min
\sum_{t=1}^T
\left[
c_t^f + c_t y_t + \overline{U}u_t
+ \underline{U}(y_{t-1} - y_t + u_t)
- \xi_t x_t
\right]
\tag{1a}
\end{equation}

\begin{equation}
\sum_{i=\max\{1,t-L+1\}}^{t} u_i \le y_t,
\qquad \forall t \in \mathcal{T}
\tag{1b}
\end{equation}

\begin{equation}
\sum_{i=\max\{1,t-\ell+1\}}^{t} u_i \le 1-\hat{y}_{t-\ell},
\qquad \forall t \in \mathcal{T}
\tag{1c}
\end{equation}

\begin{equation}
\hat{y}_{t-\ell} =
\begin{cases}
y_{t-\ell}, & t > \ell,\\
y_0, & t \le \ell,
\end{cases}
\qquad \forall t \in \mathcal{T}
\tag{1c'}
\end{equation}

\begin{equation}
y_t - y_{t-1} - u_t \le 0,
\qquad \forall t \in \mathcal{T}
\tag{1d}
\end{equation}

\begin{equation}
\underline{C} y_t \le x_t \le \overline{C} y_t,
\qquad \forall t \in \mathcal{T}
\tag{1e}
\end{equation}

\begin{equation}
x_t - x_{t-1} \le V y_{t-1} + \overline{V}(1-y_{t-1}),
\qquad \forall t \in \mathcal{T}
\tag{1f}
\end{equation}

\begin{equation}
x_{t-1} - x_t \le V y_t + \overline{V}(1-y_t),
\qquad \forall t \in \mathcal{T}
\tag{1g}
\end{equation}

\begin{equation}
\sum_{t=1}^T x_t \le Q
\tag{1h}
\end{equation}

\begin{equation}
c_t^f = \Delta t \cdot \operatorname{PWL}_t(x_t;(\chi_{tm},\phi_{tm})_{m\in\mathcal{M}_t}),
\qquad
\forall t \in \mathcal{T}
\tag{1i}
\end{equation}

\begin{equation}
y_0 = \bar{y}_0,\quad x_0 = \bar{x}_0
\tag{1j}
\end{equation}

\begin{equation}
y_t \in \{0,1\},\quad
u_t \in \{0,1\},\quad
x_t \ge 0,\quad
c_t^f \ge 0,
\qquad \forall t \in \mathcal{T}
\tag{1k}
\end{equation}

\end{document}
